% ========================================
% Chapter 4: Examples of One-Period Models
% ========================================
\chapter{Examples of One-Period Models}

\section{The Binomial Model}
We consider a one-period market $T=1$ with a universe $\Omega = \{\omega_u, \omega_d\}$. The historical probability is $P(\omega_u)=p > 0$ and $P(\omega_d)=1-p$.
\begin{itemize}
    \item A risk-free asset whose discounted price is constant: $S_0^0 = 1, S_1^0 = e^r$. The interest rate is $r > 0$.
    \item A risky asset with initial price $S_0^1=s$ and final prices $S_1^1(\omega_u) = su$ and $S_1^1(\omega_d) = sd$, with $u > d$.
\end{itemize}

\begin{proposition}
The no-arbitrage assumption (NA) is satisfied if and only if $d < e^r < u$.
In this case, there exists a unique equivalent martingale probability measure $Q \in \mathcal{M}^e(S)$, defined by:
\begin{align*}
    q = Q(\omega_u) = \frac{e^r - d}{u - d} \\
    1-q = Q(\omega_d) = \frac{u - e^r}{u - d}
\end{align*}
Since $\mathcal{M}^e(S)$ is a singleton, the market is complete.
\end{proposition}

\begin{proof}
(NA) $\iff \mathcal{M}^e(S) \neq \emptyset$. We seek $Q$ equivalent to $P$ such that $E_Q[\hat{S}_1^1] = \hat{S}_0^1$, where $\hat{S}$ is the discounted price.
\begin{equation}
E_Q\left[\frac{S_1^1}{S_1^0}\right] = \frac{S_0^1}{S_0^0} \iff q \frac{su}{e^r} + (1-q) \frac{sd}{e^r} = s
\end{equation}
Simplifying by $s$ and multiplying by $e^r$, we obtain:
\begin{equation}
q u + (1-q) d = e^r \iff q(u-d) = e^r - d
\end{equation}
This equation has a unique solution $q = \frac{e^r - d}{u - d}$. The equivalence conditions $q>0$ and $1-q>0$ are equivalent to $e^r-d>0$ and $u-e^r>0$, i.e., $d < e^r < u$.
\end{proof}

\subsection{Pricing and Hedging a Contingent Asset}
Consider a contingent asset $\hat{B}$ whose payoff at maturity is $\hat{B}(\omega_u) = \hat{B}_u$ and $\hat{B}(\omega_d) = \hat{B}_d$.
Since the market is complete, there exists a unique no-arbitrage price $A_0$ and a unique replication portfolio $(H_0^0, H_0^1)$.

\paragraph{No-Arbitrage Price:}
The price is given by the expectation of the discounted payoff under the martingale measure $Q$:
\begin{equation}
A_0 = E_Q\left[\frac{\hat{B}}{e^r}\right] = \frac{1}{e^r} \left( q \hat{B}_u + (1-q) \hat{B}_d \right) \quad \text{with} \quad q = \frac{e^r - d}{u - d}.
\end{equation}

\paragraph{Hedging Portfolio:}
The portfolio $(H_0^0, H_0^1)$ must satisfy $A_0 = H_0^0 S_0^0 + H_0^1 S_0^1 = H_0^0 + H_0^1 s$.
The portfolio value at time $T=1$ must replicate the payoff:
\begin{equation}
H_0^0 S_1^0 + H_0^1 S_1^1 = \hat{B}
\end{equation}
This leads to the system of equations:
\begin{equation}
\begin{cases}
    H_0^0 e^r + H_0^1 su = \hat{B}_u \\
    H_0^0 e^r + H_0^1 sd = \hat{B}_d
\end{cases}
\end{equation}
Subtracting the two lines gives $H_0^1 s(u-d) = \hat{B}_u - \hat{B}_d$. We deduce the quantity of risky asset to hold (called the "Delta"):
\begin{equation}
H_0^1 = \frac{\hat{B}_u - \hat{B}_d}{s(u-d)}
\end{equation}
And the quantity of risk-free asset:
\begin{equation}
H_0^0 = A_0 - H_0^1 s
\end{equation}
In the binomial model with $d < e^r < u$, I can explicitly price and hedge any contingent asset.

\section{The Trinomial Model}
We extend the previous model to a three-state universe $\Omega = \{\omega_u, \omega_m, \omega_d\}$. The final prices of the risky asset are $su, sm, sd$ with $u > m > d$.

\begin{proposition}
The no-arbitrage assumption (NA) is satisfied if and only if $d < e^r < u$.
\end{proposition}

\begin{proof}

(NA) $\iff \mathcal{M}^e(S) \neq \emptyset$. We seek a probability $Q=(q_1, q_2, q_3)$ with $q_i = Q(\omega_i) > 0$ such that:
\begin{align}
    q_1 + q_2 + q_3 &= 1 \quad (P_1) \text{ (probability condition)} \\
    q_1 u + q_2 m + q_3 d &= e^r \quad (P_2) \text{ (martingale condition)}
\end{align}

\begin{tcolorbox}[title=Geometric Visualization of Arbitrage, colback=yellow!10!white, colframe=yellow!50!black]
The set of probabilities is represented by a \textbf{simplex} (red triangle). For there to be no arbitrage (NA), the martingale plane $(P_2)$ must cross the interior of this triangle. This occurs if and only if $d < e^r < u$.
\end{tcolorbox}
\end{proof}

\paragraph{How to Read Figure 4.1:}
This figure illustrates the structure of the set of equivalent martingale measures $\mathcal{M}^e(S)$. Several key elements should be observed:
\begin{itemize}
    \item \textbf{The Simplex (Red):} It contains all probabilities $Q$ such that $\sum q_i = 1$. The interior represents \textit{equivalent} probabilities ($q_i > 0$).
    \item \textbf{The Martingale Set ($\mathcal{M}^e(S)$):} This is the line segment (blue or green) resulting from the intersection of the simplex with the martingale plane.
    \item \textbf{The Dashed Square (Gray):} On the horizontal plane ($q_1, q_3$), it delimits the domain $[0,1] \times [0,1]$. It represents the space where each individual probability is valid. The simplex only occupies the "lower" part ($q_1+q_3 \le 1$) because the sum of the three probabilities must be exactly 1.
    \item \textbf{Dependence on Parameters:} Depending on whether the "middle" return $m$ is greater or less than the risk-free rate $e^r$, the segment pivots to connect different axes of the simplex.
\end{itemize}

\begin{figure}[H]
\centering
\tdplotsetmaincoords{70}{120}
\begin{tikzpicture}[tdplot_main_coords, scale=4]
    \coordinate (O) at (0,0,0);
    
    % Axes setup:
    % q2 Up -> Z axis (0,0,1)
    % q1 Right -> Y axis (0,1,0)
    % q3 Down/Left -> X axis (1,0,0)
    
    \draw[-{Latex[length=2mm]}] (O) -- (1.5,0,0) node[anchor=north east]{$q_3$};
    \draw[-{Latex[length=2mm]}] (O) -- (0,1.5,0) node[anchor=north west]{$q_1$};
    \draw[-{Latex[length=2mm]}] (O) -- (0,0,1.5) node[anchor=south]{$q_2$};
    
    % Unit Square on floor (q1, q3)
    \draw[dashed, gray!80] (1,0,0) -- (1,1,0) -- (0,1,0);
    \node[gray, font=\tiny, anchor=south west] at (1,1,0) {$(1,1)$};

    % Vertices of Simplex on axes (adjusted to match professor's schema)
    \coordinate (Q3) at (1.25,0,0); % q3 axis intercept - outside unit square
    \coordinate (Q1) at (0,0.75,0); % q1 axis intercept - inside unit square
    \coordinate (Q2) at (0,0,1); % q2 axis intercept

    % Simplex (P1) - RED
    \draw[fill=red!5, opacity=0.6] (Q1) -- (Q2) -- (Q3) -- cycle;
    \draw[red, thick] (Q1) -- (Q2);
    \draw[red, thick] (Q2) -- (Q3);
    \draw[red, thick] (Q3) -- (Q1);
    
    \node[red, font=\bfseries] at (0.2, 0.2, 0.9) {Simplex $(P_1)$};

    % Blue Setup: Line on Q1-Q3 edge to Q2-Q3 edge
    % Point on q1-q3 edge (floor): parametric from (0,0.75,0) to (1.25,0,0)
    \coordinate (Blue_P1) at (0.9, 0.2, 0); 
    % Point on q2-q3 edge (x-z plane): parametric from (0,0,1) to (1.25,0,0)
    \coordinate (Blue_P2) at (0.85, 0, 0.3);
    
    \draw[ultra thick, blue] (Blue_P1) -- (Blue_P2) node[pos=0.1, left, xshift=-10pt, yshift=5pt, font=\scriptsize]{$\mathcal{M}^e(S)$};
    \fill[blue] (Blue_P1) circle (0.4pt);
    \fill[blue] (Blue_P2) circle (0.4pt);
    
    % Blue Normal n1
    \draw[->, blue, thick] (O) -- (1.5, 1.5, 1.5) node[above right, blue, inner sep=2pt] {$\vec{n}_1$};

    % Green Setup: Line on Q1-Q3 edge to Q1-Q2 edge
    % Point on q1-q3 edge (floor): closer to Q1
    \coordinate (Green_P1) at (0.4, 0.5, 0);
    % Point on q1-q2 edge (y-z plane): parametric from (0,0.75,0) to (0,0,1)
    \coordinate (Green_P2) at (0, 0.4, 0.45);

    \draw[ultra thick, green!60!black] (Green_P1) -- (Green_P2) node[pos=0.1, right, xshift=10pt, yshift=5pt, font=\scriptsize]{$\mathcal{M}^e(S)$};
    \fill[green!60!black] (Green_P1) circle (0.4pt);
    \fill[green!60!black] (Green_P2) circle (0.4pt);
    
    % Green Normal n2
    \draw[->, green!60!black, thick] (O) -- (0.2, 0.9, 0.4) node[right, green!60!black, inner sep=3pt] {$\vec{n}_2$};

    % Labels for points
    \node[below=8pt, blue] at (Blue_P1) {$\bar{P}_1$}; 
    \node[left=10pt, blue] at (Blue_P2) {$\bar{P}_2$};
    
    \node[below right=8pt, green!60!black] at (Green_P1) {$\bar{P}_1$}; 
    \node[right=10pt, green!60!black] at (Green_P2) {$\bar{P}_2'$};

\end{tikzpicture}
\caption{Geometry of incompleteness in the trinomial model. The intersection (thick segment) between the probability space and the martingale constraint shows the existence of infinitely many RNM measures. The blue configuration corresponds to $m \ge e^r$ and the green to $m < e^r$.}
\end{figure}

\begin{tcolorbox}[title=The Interest of the Geometric View, colback=red!5!white, colframe=red!50!black]
This view allows understanding that the price of an asset is no longer unique: it can vary along the segment $\mathcal{M}^e(S)$. The no-arbitrage price bounds are simply the asset values at the segment endpoints ($\bar{P}_1$ and $\bar{P}_2$).
\end{tcolorbox}

More generally, to obtain $\mathcal{M}^e(S)$, we seek the intersection between the interior of the simplex (open) and the plane $(P_2)$ with equation:
\begin{equation}
    q_1 u + q_2 m + q_3 d = e^r
\end{equation}
The following vector is orthogonal to plane $(P_2)$:
\begin{equation}
    \vec{V} = \begin{pmatrix} u/e^r \\ m/e^r \\ d/e^r \end{pmatrix}
\end{equation}
Under the no-arbitrage assumption (NA), we always have:
\begin{equation}
    \frac{u}{e^r} > 1 \quad \text{and} \quad 0 < \frac{d}{e^r} < 1
\end{equation}
The intersection then depends on the position of $\frac{m}{e^r}$ relative to 1. We distinguish two cases (illustrated in Figure 4.1):
To find the price bounds of a claim, it suffices to evaluate its expectation at the two extreme points of this segment, denoted $\bar{P}_1$ and $\bar{P}_2$.

The exact form of these extreme points depends on the position of the intermediate return $m$ relative to the risk-free rate $e^r$.
The set of martingale measures $\mathcal{M}^e(S)$ is the intersection of two geometric constraints: the probability simplex ($q_1+q_2+q_3=1, q_i>0$) and the martingale condition plane ($q_1u + q_2m + q_3d = e^r$). In the trinomial case with $d < e^r < u$, this intersection is a line segment. The no-arbitrage price bounds are attained at the endpoints of this segment. These extreme points correspond to martingale measures that "ignore" one of the three possible states, i.e., by assigning it zero probability ($q_i=0$).

Let us detail the calculation for the case where $m \ge e^r$.

\begin{itemize}
    \item \textbf{Calculation of $\bar{P}_1$:} This point is obtained by assuming that the intermediate state $m$ has zero probability, i.e., $q_2=0$. The system of equations reduces to:
    \begin{enumerate}
        \item $q_1 + q_3 = 1$
        \item $q_1u + q_3d = e^r$
    \end{enumerate}
    Substituting $q_3 = 1 - q_1$ into the second equation, we get $q_1u + (1 - q_1)d = e^r$, which gives $q_1(u - d) = e^r - d$. The solution is:
    \[
    q_1 = \frac{e^r - d}{u - d} \quad \text{and} \quad q_3 = 1 - q_1 = \frac{u - e^r}{u - d}.
    \]
    Thus $\bar{P}_1 = ( \frac{e^r-d}{u-d}, 0, \frac{u-e^r}{u-d} )$.

    \item \textbf{Calculation of $\bar{P}_2$:} This point is obtained by setting $q_1=0$. The system becomes:
    \begin{enumerate}
        \item $q_2 + q_3 = 1$
        \item $q_2m + q_3d = e^r$
    \end{enumerate}
    The resolution is analogous and gives:
    \[
    q_2 = \frac{e^r - d}{m - d} \quad \text{and} \quad q_3 = \frac{m - e^r}{m - d}.
    \]
    Thus $\bar{P}_2 = ( 0, \frac{e^r-d}{m-d}, \frac{m-e^r}{m-d} )$.
\end{itemize}

The calculation for the case $m < e^r$ follows the same logic, setting $q_2=0$ then $q_3=0$.

\subsection{Pricing of Contingent Assets}
For a contingent asset $\hat{B}$ with payoffs $(\hat{B}_u, \hat{B}_m, \hat{B}_d)$, the set of no-arbitrage prices is the open interval $]\underline{\Pi}(\hat{B}), \overline{\Pi}(\hat{B})[$, where:
\begin{itemize}
    \item $\underline{\Pi}(\hat{B}) = \inf_{Q \in \mathcal{M}^e(S)} E_Q[\hat{B}/e^r]$
    \item $\overline{\Pi}(\hat{B}) = \sup_{Q \in \mathcal{M}^e(S)} E_Q[\hat{B}/e^r]$
\end{itemize}
The map $Q \mapsto E_Q[\cdot]$ being linear and the set $\mathcal{M}^e(S)$ being convex (a segment), the supremum and infimum are attained at the extreme points of the segment. These extreme points, denoted $\bar{P}_1$ and $\bar{P}_2$, correspond to cases where one of the states has zero probability.

Two cases arise depending on the position of $m$ relative to $e^r$:
\begin{enumerate}
    \item If $m \ge e^r$: The extreme points are $\bar{P}_1$ (intersection with $q_2=0$) and $\bar{P}_2$ (intersection with $q_1=0$).
        \begin{itemize}
            \item $\bar{P}_1 = (\frac{e^r-d}{u-d}, 0, \frac{u-e^r}{u-d})$
            \item $\bar{P}_2 = (0, \frac{e^r-d}{m-d}, \frac{m-e^r}{m-d})$
        \end{itemize}
    \item If $m < e^r$: The extreme points are $\bar{P}_1$ (intersection with $q_2=0$) and $\bar{P}_2$ (intersection with $q_3=0$).
        \begin{itemize}
            \item $\bar{P}_1 = (\frac{e^r-d}{u-d}, 0, \frac{u-e^r}{u-d})$
            \item $\bar{P}_2 = (\frac{e^r-m}{u-m}, \frac{u-e^r}{u-m}, 0)$
        \end{itemize}
\end{enumerate}

\begin{example}[European Call Option]
Consider a call option with strike $K$ on the risky asset, so that the discounted payoff is $\hat{B} = (S_1^1 - K)^+/e^r$. Suppose $su > K > sm$ and $sd < K$. The payoffs are:
\[
\hat{B}(\omega_u) = \frac{su-K}{e^r}, \quad \hat{B}(\omega_m) = 0, \quad \hat{B}(\omega_d) = 0
\]
The price bounds are then:
\begin{itemize}
    \item $\overline{\Pi}(\hat{B}) = \max(E_{\bar{P}_1}[\hat{B}/e^r], E_{\bar{P}_2}[\hat{B}/e^r])$.
    Since $u$ is the only state where the payoff is non-zero, the sup is attained for the probability that maximizes $q_1$. This is $\bar{P}_1$ if $m < e^r$ or $\bar{P}_2$ if $m > e^r$ (depending on values).
    Suppose $m < e^r$. The max of $q_1$ is at $\bar{P}_2$.
    $\overline{\Pi}(\hat{B}) = E_{\bar{P}_2}[\hat{B}/e^r] = \frac{e^r-m}{u-m} \cdot \frac{su-K}{e^r}$.
    \item $\underline{\Pi}(\hat{B}) = \min(E_{\bar{P}_1}[\hat{B}/e^r], E_{\bar{P}_2}[\hat{B}/e^r])$.
    The min is at $\bar{P}_1$.
    $\underline{\Pi}(\hat{B}) = E_{\bar{P}_1}[\hat{B}/e^r] = \frac{e^r-d}{u-d} \cdot \frac{su-K}{e^r}$.
\end{itemize}
\end{example}

\begin{example}[Numerical Exercise - Trinomial Model]
Consider a one-period trinomial model with the following parameters: risk-free rate $r=0$ (hence $e^r = 1$), initial price $S_0^1 = 1$, return factors $u=1.2$, $m=1$, $d=0.8$. We consider a call option with strike $K = 0.9$, whose discounted payoff is $\hat{B} = (\hat{S}^1_1 - 0.9)^+$.

\textbf{Questions:}
\begin{enumerate}
    \item Verify the no-arbitrage condition.
    \item Calculate the payoffs in each state.
    \item Determine the extreme martingale measures.
    \item Deduce the no-arbitrage price interval.
    \item Construct a super-replication portfolio.
\end{enumerate}
\end{example}

\section*{Solution to Numerical Exercise 4.3}

This subsection provides a detailed solution to the above exercise.

\begin{enumerate}
    \item \textbf{Parameters:}
    \begin{itemize}
        \item Risk-free rate: $r=0$, hence $e^r = 1$.
        \item Initial price: $S_0^1 = 1$.
        \item Return factors: $u=1.2, m=1, d=0.8$.
        \item Contingent claim: $\hat{B} = (\hat{S}^1_1 - 0.9)^+$.
    \end{itemize}
    \item \textbf{Market Analysis:} We verify the no-arbitrage condition: $d < e^r < u$ becomes $0.8 < 1 < 1.2$. The condition is satisfied. The market is incomplete.
    \item \textbf{Payoff Calculation:}
    \begin{itemize}
        \item State $u$: $\hat{S}^1_1 = 1.2$, $\hat{B}_u = (1.2 - 0.9)^+ = 0.3$.
        \item State $m$: $\hat{S}^1_1 = 1.0$, $\hat{B}_m = (1.0 - 0.9)^+ = 0.1$.
        \item State $d$: $\hat{S}^1_1 = 0.8$, $\hat{B}_d = (0.8 - 0.9)^+ = 0$.
    \end{itemize}
    \item \textbf{Determination of Extreme Measures:} We are in the case $m = e^r = 1$, which is a limiting case of $m \ge e^r$.
    \begin{itemize}
        \item $\bar{P}_1$ (with $q_2=0$, hence $q_1+q_3=1$ and $1.2q_1 + 0.8q_3 = 1$):
        \[
        q_1 = \frac{1-0.8}{1.2-0.8} = \frac{0.2}{0.4} = 0.5.
        \]
        Thus $\bar{P}_1 = (0.5, 0, 0.5)$.
        \item $\bar{P}_2$ (with $q_1=0$, hence $q_2+q_3=1$ and $1.0q_2 + 0.8q_3 = 1$):
        \[
        q_2 = \frac{1-0.8}{1-0.8} = 1.
        \]
        Thus $\bar{P}_2 = (0, 1, 0)$.
    \end{itemize}
    \item \textbf{Price Bound Calculation:} The no-arbitrage price interval is $]\underline{\Pi}(\hat{B}), \overline{\Pi}(\hat{B})[$:
    \begin{itemize}
        \item $E_{\bar{P}_1}[\hat{B}] = 0.5 \times 0.3 + 0 \times 0.1 + 0.5 \times 0 = 0.15$.
        \item $E_{\bar{P}_2}[\hat{B}] = 0 \times 0.3 + 1 \times 0.1 + 0 \times 0 = 0.1$.
    \end{itemize}
    The price interval is therefore $]0.1, 0.15[$.
    \item \textbf{Super-Replication Portfolio Calculation:} The minimum cost to super-replicate is $\alpha = \overline{\Pi}(\hat{B}) = 0.15$. We need to find a portfolio $(H_0^0, H_0^1)$ such that its initial cost is 0.15 and its final value dominates the payoff $\hat{B}$ in all states.
    \begin{itemize}
        \item Initial cost: $\hat{V}_0 = H_0^0 \cdot 1 + H_0^1 \cdot 1 = 0.15$.
        \item The proposed portfolio is $H_0^1 = 0.75$ and $H_0^0 = -0.6$.
        \item Cost verification: $-0.6 + 0.75 = 0.15$. Correct.
        \item Super-replication verification at $T=1$:
        \begin{itemize}
            \item State $u$: $\hat{V}_1 = -0.6 \times 1 + 0.75 \times 1.2 = -0.6 + 0.9 = 0.3$. Now $\hat{B}_u = 0.3$. So $\hat{V}_1 \ge \hat{B}_u$. OK.
            \item State $m$: $\hat{V}_1 = -0.6 \times 1 + 0.75 \times 1.0 = 0.15$. Now $\hat{B}_m = 0.1$. So $\hat{V}_1 \ge \hat{B}_m$ ($0.15 \ge 0.1$). OK.
            \item State $d$: $\hat{V}_1 = -0.6 \times 1 + 0.75 \times 0.8 = -0.6 + 0.6 = 0$. Now $\hat{B}_d = 0$. So $\hat{V}_1 \ge \hat{B}_d$. OK.
        \end{itemize}
    \end{itemize}
    The portfolio successfully super-replicates the payoff for a cost of 0.15.
\end{enumerate}
